\section{Discussion and Limitations}\label{sec:discussion}
The studies support a common practical motivation: intermediate feedback can distinguish useful partial behavior when terminal verification supplies little contrast. Their evidential roles differ. LIBERO provides a reward comparison within one reported small-compute setup; the CDPR evidence shows continued improvement of a staged residual controller. Different backbones, action likelihoods, tasks, and compute budgets prevent interpreting the two as a controlled comparison of dense and sparse feedback.

Stage awareness also leaves objective mismatch unresolved. A reward can favor repeated partial competence over rare completion, and protecting scalar prefix credit cannot eliminate interference in a shared network. Conditional latent scoring fixes the recorded sampling model but does not certify the full selected, normalized, and clipped optimizer. The finite-model analysis should therefore be interpreted alongside measured learning outcomes, with neural-policy alignment remaining an open diagnostic.

The CDPR results reuse development scenes and one training lineage. Independent seeds, a locked final test, and matched credit/likelihood/prior ablations are needed for stronger attribution and generalization claims. LIBERO lacks trial denominators and quantitative throughput records in the supplied draft. The physical success narrative is a writing assumption until trial evidence is inserted. Historical scripted initialization, manual rewards, calibration, inherited compute, and limited object/destination coverage further constrain claims of autonomous adaptation. Action-prior ablations are especially important because pretrained visual features and pretrained action knowledge can contribute differently.
